\documentclass[12pt,a4paper]{article}

% Documento di notazione per Scientific Workplace 5.5.
% Si usano pacchetti standard e una sintassi LaTeX prudente.
\usepackage[T1]{fontenc}
\usepackage{amsmath,amssymb,amsthm}
\usepackage{array}
\usepackage{longtable}
\usepackage{geometry}
\geometry{margin=2.7cm}

\newcommand{\R}{\mathbb{R}}
\newcommand{\N}{\mathbb{N}}
\newcommand{\E}{\mathbb{E}}
\newcommand{\Prob}{\mathbb{P}}
\newcommand{\Q}{\mathbb{Q}}
\newcommand{\F}{\mathcal{F}}
\newcommand{\B}{\mathcal{B}}
\newcommand{\G}{\mathcal{G}}
\newcommand{\1}{\mathbf{1}}
\newcommand{\Var}{\operatorname{Var}}
\newcommand{\Cov}{\operatorname{Cov}}
\newcommand{\VaR}{\operatorname{VaR}}
\newcommand{\CVaR}{\operatorname{CVaR}}
\newcommand{\argmin}{\operatorname*{arg\,min}}
\newcommand{\argmax}{\operatorname*{arg\,max}}
\newcommand{\EVPI}{\operatorname{EVPI}}
\newcommand{\VSS}{\operatorname{VSS}}
\setlength{\headheight}{14pt}

\title{MQF\_Notazione\\[0.2cm]
\large Convenzioni di notazione per il corso\\
\large Metodi Quantitativi per la Finanza}
\author{}
\date{Versione preliminare}

\begin{document}
\maketitle

\section{Finalit\`{a} del documento}

Questo documento raccoglie le convenzioni di notazione da utilizzare nel manuale, nelle slides, negli esercizi e nelle applicazioni Python del corso \emph{Metodi Quantitativi per la Finanza}. La finalit\`{a} principale \`{e} garantire coerenza formale tra le diverse parti del materiale didattico.

Il corso lavora prevalentemente in tempo discreto. Salvo diversa indicazione, l'indice temporale \`{e}
\[
 t=0,1,\ldots,T,
\]
dove $T$ indica l'orizzonte finale del modello. I simboli introdotti in questo documento devono essere considerati convenzioni generali. Eventuali eccezioni locali devono essere esplicitate nel testo della singola lezione.

\section{Convenzioni tipografiche generali}

\begin{longtable}{>{\raggedright}p{3.2cm}>{\raggedright}p{4.2cm}>{\raggedright\arraybackslash}p{7.0cm}}
\hline
\textbf{Oggetto} & \textbf{Notazione} & \textbf{Uso previsto}\\
\hline
Scalari & $a,b,c,\alpha,\beta$ & Parametri reali, coefficienti, livelli di confidenza, tassi o pesi.\\
Vettori & $x,y,\lambda,\xi$ & Vettori colonna, salvo diversa indicazione:
$x$ per decisioni, $y_s$ per decisioni di secondo stadio, $\lambda$ per
variabili duali, $\xi$ per dati aleatori.\\
Matrici & $A,B,P,Q$ & Matrici di coefficienti, matrici di transizione,
matrici di vincoli.\\
Insiemi & $A,B,C,\Omega,\mathcal{S},\mathcal{I}$ & Eventi, spazio campionario,
insieme degli scenari, insieme degli stati markoviani.\\
Variabili casuali & $X,Y,Z,L,H$ & Variabili casuali reali; $L$ per perdite,
$H$ per payoff terminali.\\
Processi & $\{X_t\}_{t=0}^{T}$, $\{S_t\}_{t=0}^{T}$ & Sequenze temporali
aleatorie in tempo discreto.\\
Valore ottimo & $z^*$ & Valore ottimo di un problema di ottimizzazione.\\
Soluzione ottima & $x^*$ & Vettore ottimo delle variabili decisionali.\\
\hline
\end{longtable}

Per i vettori si usa, quando necessario, la notazione trasposta $x'$. Ad esempio, il prodotto scalare tra $c$ e $x$ pu\`{o} essere scritto come
\[
 c'x.
\]
In alternativa, quando la chiarezza lo richiede, si pu\`{o} usare $c^{\top}x$. Nei materiali del corso si preferisce $c'x$ per semplicit\`{a} tipografica.

\section{Spazi di probabilit\`{a} ed eventi}

Lo spazio di probabilit\`{a} di riferimento \`{e}
\[
 (\Omega,\F,\Prob),
\]
dove $\Omega$ \`{e} lo spazio degli esiti elementari, $\F$ \`{e} una sigma-algebra di eventi e $\Prob$ \`{e} una misura di probabilit\`{a}.

\begin{longtable}{>{\raggedright}p{3.4cm}>{\raggedright}p{4.4cm}>{\raggedright\arraybackslash}p{6.6cm}}
\hline
\textbf{Concetto} & \textbf{Notazione} & \textbf{Significato}\\
\hline
Esito elementare & $\omega\in\Omega$ & Stato elementare dell'incertezza.\\
Evento & $A\in\F$ & Sottoinsieme misurabile di $\Omega$.\\
Complementare & $A^c$ & Evento complementare di $A$.\\
Unione & $A\cup B$ & Evento ``$A$ oppure $B$''.\\
Intersezione & $A\cap B$ & Evento ``$A$ e $B$''.\\
Probabilit\`{a} & $\Prob(A)$ & Probabilit\`{a} dell'evento $A$.\\
Probabilit\`{a} condizionata & $\Prob(A\mid B)$ & Probabilit\`{a} di $A$ condizionatamente a $B$.\\
\hline
\end{longtable}

Per eventi $A,B\in\F$, con $\Prob(B)>0$, la probabilit\`{a} condizionata \`{e}
\[
 \Prob(A\mid B)=\frac{\Prob(A\cap B)}{\Prob(B)}.
\]
La regola del prodotto \`{e}
\[
 \Prob(A\cap B)=\Prob(A\mid B)\Prob(B).
\]
Se $\{B_1,\ldots,B_n\}$ \`{e} una partizione di $\Omega$, con $\Prob(B_i)>0$, allora
\[
 \Prob(A)=\sum_{i=1}^{n}\Prob(A\mid B_i)\Prob(B_i).
\]
La formula di Bayes \`{e}
\[
 \Prob(B_j\mid A)=\frac{\Prob(A\mid B_j)\Prob(B_j)}{\sum_{i=1}^{n}\Prob(A\mid B_i)\Prob(B_i)}.
\]

\section{Variabili casuali, distribuzioni e momenti}

Una variabile casuale reale \`{e} indicata con una lettera maiuscola, ad esempio
\[
 X:\Omega\rightarrow \R.
\]
Le realizzazioni sono indicate con lettere minuscole, ad esempio $x$.

\begin{longtable}{>{\raggedright}p{3.8cm}>{\raggedright}p{4.2cm}>{\raggedright\arraybackslash}p{6.4cm}}
\hline
\textbf{Concetto} & \textbf{Notazione} & \textbf{Significato}\\
\hline
Variabile casuale & $X$ & Quantit\`{a} aleatoria, ad esempio rendimento, perdita o payoff.\\
Realizzazione & $x$ & Valore assunto da $X$.\\
Funzione di ripartizione & $F_X(x)$ & $F_X(x)=\Prob(X\leq x)$.\\
Funzione di massa & $p_X(x)$ & $p_X(x)=\Prob(X=x)$ per variabili discrete.\\
Densit\`{a} & $f_X(x)$ & Densit\`{a} di una variabile continua.\\
Valore atteso & $\E[X]$ & Media teorica di $X$.\\
Varianza & $\Var(X)$ & Misura di dispersione di $X$.\\
Covarianza & $\Cov(X,Y)$ & Dipendenza lineare tra $X$ e $Y$.\\
Indicatrice & $\1_A$ & Variabile che vale $1$ su $A$ e $0$ altrove.\\
\hline
\end{longtable}

Per una variabile casuale discreta,
\[
 \E[X]=\sum_x x\,p_X(x).
\]
Per una variabile casuale continua con densit\`{a} $f_X$,
\[
 \E[X]=\int_{-\infty}^{+\infty}x f_X(x)\,dx.
\]
La varianza \`{e}
\[
 \Var(X)=\E\left[(X-\E[X])^2\right]=\E[X^2]-\left(\E[X]\right)^2.
\]
La covarianza tra $X$ e $Y$ \`{e}
\[
 \Cov(X,Y)=\E\left[(X-\E[X])(Y-\E[Y])\right].
\]

\section{Informazione e valori attesi condizionati}

L'informazione disponibile \`{e} rappresentata mediante sigma-algebre. Se $\G\subseteq\F$, allora $\G$ rappresenta un livello informativo meno fine di $\F$.

Il valore atteso condizionato di $X$ rispetto a $\G$ \`{e} indicato con
\[
 \E[X\mid \G].
\]
Quando il condizionamento \`{e} rispetto a un evento $A$ con $\Prob(A)>0$, si scrive
\[
 \E[X\mid A].
\]

Se $X$ \`{e} discreta,
\[
 \E[X\mid A]=\sum_x x\,\Prob(X=x\mid A).
\]
Se $\{B_1,\ldots,B_n\}$ \`{e} una partizione informativa di $\Omega$, allora
\[
 \E[X\mid \G]=\sum_{i=1}^{n}\E[X\mid B_i]\1_{B_i},
\]
dove $\G=\sigma(B_1,\ldots,B_n)$.

Le propriet\`{a} fondamentali da usare nel corso sono:
\[
 \E[aX+bY\mid\G]=a\E[X\mid\G]+b\E[Y\mid\G],
\]
\[
 \E\left[\E[X\mid\G]\right]=\E[X],
\]
\[
 \E\left[\E[X\mid\F_t]\mid\F_s\right]=\E[X\mid\F_s],\qquad s\leq t.
\]
L'ultima identit\`{a} \`{e} la propriet\`{a} iterata del valore atteso condizionato, o \emph{tower property}.

\section{Tempo discreto, processi stocastici e filtrazioni}

Un processo stocastico in tempo discreto \`{e} indicato come
\[
 \{X_t\}_{t=0}^{T}.
\]
Quando il processo rappresenta il prezzo di un'attivit\`{a} finanziaria, si usa preferibilmente
\[
 \{S_t\}_{t=0}^{T}.
\]

La filtrazione \`{e}
\[
 \{\F_t\}_{t=0}^{T},
\]
dove $\F_t$ rappresenta l'informazione disponibile al tempo $t$. Un processo $\{X_t\}_{t=0}^{T}$ \`{e} adattato alla filtrazione se $X_t$ \`{e} $\F_t$-misurabile per ogni $t$.

Una traiettoria campionaria \`{e} la sequenza
\[
 X_0(\omega),X_1(\omega),\ldots,X_T(\omega)
\]
associata a uno specifico esito $\omega\in\Omega$.

\section{Scenari}

Nei modelli discreti di decisione sotto incertezza, l'insieme degli scenari \`{e} indicato con
\[
 \mathcal{S}=\{1,\ldots,N\}.
\]
Lo scenario generico \`{e} indicato con $s\in\mathcal{S}$ e la sua probabilit\`{a} con
\[
 p_s,
 \qquad p_s\geq 0,
 \qquad \sum_{s\in\mathcal{S}}p_s=1.
\]
Un parametro dipendente dallo scenario pu\`{o} essere scritto come $a_s$, $b_s$, $c_s$ oppure, in forma vettoriale, come $\xi_s$.

Il vettore aleatorio sottostante \`{e} indicato con
\[
 \xi.
\]
Quando si distinguono scenari discreti del modello e repliche simulate, si
riserva $N$ al numero di scenari e si usa $M$ per il numero di repliche
simulate. L'indice di replica simulata \`{e} indicato con
\[
 k=1,\ldots,M.
\]
Una realizzazione simulata della variabile casuale $X$ nella replica $k$ pu\`{o}
essere indicata con
\[
 X^{(k)}.
\]
Questa convenzione \`{e} usata soprattutto nelle applicazioni Python e nelle
stime Monte Carlo.

\section{Catene di Markov}

Per le catene di Markov in tempo discreto, l'insieme degli stati \`{e} indicato con
\[
 \mathcal{I}=\{1,\ldots,m\}.
\]
La catena \`{e}
\[
 \{X_t\}_{t=0}^{T},
\]
dove $X_t\in\mathcal{I}$. La matrice di transizione \`{e}
\[
 P=(p_{ij})_{i,j=1}^{m},
\]
con
\[
 p_{ij}=\Prob(X_{t+1}=j\mid X_t=i),
 \qquad
 p_{ij}\geq 0,
 \qquad
 \sum_{j=1}^{m}p_{ij}=1.
\]

La propriet\`{a} di Markov \`{e}
\[
 \Prob(X_{t+1}=j\mid X_t=i,X_{t-1}=i_{t-1},\ldots,X_0=i_0)
 =\Prob(X_{t+1}=j\mid X_t=i).
\]

Se $\pi_t$ indica il vettore riga delle probabilit\`{a} di stato al tempo $t$, allora
\[
 \pi_{t+1}=\pi_t P,
 \qquad
 \pi_t=\pi_0P^t.
\]
Uno stato $i$ \`{e} assorbente se
\[
 p_{ii}=1.
\]

\section{Rendimenti, perdite e misure di rischio}

Il prezzo di un'attivit\`{a} finanziaria al tempo $t$ \`{e} indicato con $S_t$. Il rendimento semplice su un periodo pu\`{o} essere indicato come
\[
 R_{t+1}=\frac{S_{t+1}-S_t}{S_t}.
\]
Quando si lavora con perdite, la variabile casuale di perdita \`{e} indicata con
\[
 L.
\]
Per convenzione, valori pi\`{u} elevati di $L$ rappresentano esiti peggiori.

Per un livello di confidenza $\alpha\in(0,1)$, il Value-at-Risk \`{e} indicato con
\[
 \VaR_{\alpha}(L).
\]
In forma generale,
\[
 \VaR_{\alpha}(L)=\inf\{\ell\in\R:F_L(\ell)\geq \alpha\}.
\]
Il Conditional Value-at-Risk \`{e} indicato con
\[
 \CVaR_{\alpha}(L).
\]
Nel caso continuo, si pu\`{o} usare l'interpretazione
\[
 \CVaR_{\alpha}(L)=\E\left[L\mid L\geq \VaR_{\alpha}(L)\right],
\]
mentre nelle formulazioni di ottimizzazione si preferisce la rappresentazione
\[
 \CVaR_{\alpha}(L)=\min_{\eta\in\R}\left\{\eta+\frac{1}{1-\alpha}\E[(L-\eta)^+]\right\},
\]
dove
\[
(x)^+=\max\{x,0\}.
\]

Nelle formulazioni discrete del CVaR si indicano con $u_s$ le variabili
ausiliarie che rappresentano gli eccessi di perdita rispetto alla soglia
$\eta$:
\[
u_s\geq L_s(x)-\eta,
\qquad
u_s\geq0,
\qquad
s\in\mathcal S.
\]
La formulazione lineare standard \`{e}
\[
\min_{x,\eta,u}
\eta+\frac{1}{1-\alpha}\sum_{s\in\mathcal S}p_s u_s.
\]
Il simbolo $z$ resta riservato ai valori ottimi dei problemi di
ottimizzazione.

\section{Modello binomiale e pricing discreto}

Nel modello binomiale, il prezzo iniziale \`{e} $S_0>0$.
In un periodo, il prezzo pu\`{o} assumere i valori
\[
S_1^u=uS_0,
\qquad
S_1^d=dS_0,
\]
con
\[
u>d>0.
\]
Il tasso privo di rischio su un periodo \`{e} indicato con $r$, per cui
il fattore di capitalizzazione \`{e} $1+r$.
Per la valutazione di flussi certi, si indica con
\[
D_t=(1+r)^{-t}
\]
il fattore di sconto a $t$ periodi. Se $c_t$ \`{e} il flusso certo pagato al
tempo $t$, il valore iniziale di una sequenza di flussi obbligazionari certi
pu\`{o} essere scritto come
\[
V_0^{obbl}=\sum_{t=1}^{T}D_t c_t.
\]

La condizione elementare di assenza di arbitraggio \`{e}
\[
d<1+r<u.
\]
La probabilit\`{a} risk-neutral a un periodo \`{e}
\[
q=\frac{1+r-d}{u-d}.
\]
Lo strike price di un'opzione \`{e} indicato con $K$. Per un'opzione europea
call e put con scadenza $T$, i payoff terminali sono rispettivamente
\[
H^{\mathrm{call}}=(S_T-K)^+,
\qquad
H^{\mathrm{put}}=(K-S_T)^+.
\]
La misura risk-neutral \`{e} indicata con $\Q$. Il valore di un derivato con
payoff terminale $H$ \`{e} indicato con $V_t$. A un periodo,
\[
V_0=\frac{1}{1+r}\E_{\Q}[H]
=\frac{1}{1+r}\left(qH^u+(1-q)H^d\right).
\]

\section{Programmazione lineare}

Un problema di programmazione lineare viene scritto, salvo diversa
indicazione, nella forma
\[
\min_x \; c'x
\]
soggetto a
\[
Ax\geq b,
\qquad
x\geq0.
\]
Il vettore delle variabili decisionali \`{e} $x$, il valore ottimo \`{e}
$z^*$ e una soluzione ottima \`{e} $x^*$. Le variabili duali sono indicate
preferibilmente con $\lambda$. Per il problema primale precedente, il duale
associato \`{e}
\[
\max_{\lambda} \; b'\lambda
\]
soggetto a
\[
A'\lambda\leq c,
\qquad
\lambda\geq0.
\]
Il simbolo $y_s$ \`{e} riservato alle decisioni di secondo stadio nei problemi
di programmazione stocastica.

\section{Programmazione stocastica}

Nei problemi stocastici a due stadi, $x$ indica la decisione di primo stadio,
presa prima dell'osservazione dello scenario. Per ogni scenario
$s\in\mathcal S$, $y_s$ indica la decisione di secondo stadio o decisione di
recourse.

Una forma standard \`{e}
\[
\min_{x,y_s} \; c'x+\sum_{s\in\mathcal S}p_s d_s'y_s
\]
soggetta a
\[
Ax\geq b,
\qquad
W_s y_s\geq h_s-T_sx,
\qquad
y_s\geq0,
\quad s\in\mathcal S.
\]
Il valore ottimo del problema stocastico \`{e} indicato con $z^{SP}$.

Il valore wait-and-see \`{e} indicato con $z^{WS}$ e corrisponde al valore
atteso dei problemi risolti separatamente dopo l'osservazione dello scenario.

Si indichi con $x^{EV}$ la soluzione ottima del problema medio. Il valore
$z^{EV}$ non \`{e} il valore ottimo del solo problema medio, ma il valore
della funzione obiettivo stocastica ottenuto imponendo la decisione di primo
stadio $x=x^{EV}$.

Per problemi formulati in minimizzazione si usa la convenzione
\[
z^{WS}\leq z^{SP}\leq z^{EV}.
\]
Pertanto,
\[
\EVPI=z^{SP}-z^{WS}\geq0,
\qquad
\VSS=z^{EV}-z^{SP}\geq0.
\]

\section{Corrispondenza tra notazione matematica e Python}

\begin{longtable}{>{\raggedright}p{4.0cm}>{\raggedright}p{4.2cm}>{\raggedright\arraybackslash}p{5.8cm}}
\hline
\textbf{Oggetto matematico} & \textbf{Nome Python suggerito} & \textbf{Commento}\\
\hline
$S_t$ & \texttt{S\_t}, \texttt{prices} & Prezzi simulati o osservati.\\
$X_t$ & \texttt{X\_t}, \texttt{states} & Processo generico o stati markoviani.\\
$P=(p_{ij})$ & \texttt{P} & Matrice di transizione.\\
$\pi_t$ & \texttt{pi\_t} & Distribuzione sugli stati.\\
$L$ & \texttt{losses} & Vettore delle perdite.\\
$\VaR_\alpha$ & \texttt{var\_alpha} & Value-at-Risk al livello $\alpha$.\\
$\CVaR_\alpha$ & \texttt{cvar\_alpha} & Conditional Value-at-Risk.\\
$\eta$ & \texttt{eta} & Soglia nella formulazione del CVaR.\\
$u_s$ & \texttt{u[s]} & Eccesso di perdita nello scenario $s$.\\
$x$ & \texttt{x} & Variabili decisionali di primo stadio o di portafoglio.\\
$\lambda$ & \texttt{lambda\_dual} & Variabili duali nei problemi di programmazione lineare.\\
$y_s$ & \texttt{y[s]} & Variabili di secondo stadio indicizzate per scenario.\\
$p_s$ & \texttt{p\_s}, \texttt{probabilities} & Probabilit\`{a} degli scenari.\\
$d_s$ & \texttt{d[s]}, \texttt{recourse\_costs} & Costi di recourse nello scenario $s$.\\
$\xi_s$ & \texttt{xi[s]}, \texttt{scenario\_data} & Dati aleatori dello scenario.\\
$z^{SP}$ & \texttt{z\_sp} & Valore ottimo del problema stocastico.\\
$z^{WS}$ & \texttt{z\_ws} & Valore wait-and-see.\\
$x^{EV}$ & \texttt{x\_ev} & Soluzione ottima del problema medio.\\
$z^{EV}$ & \texttt{z\_ev} & Valore stocastico ottenuto fissando $x=x^{EV}$.\\
$\EVPI$ & \texttt{evpi} & Expected Value of Perfect Information.\\
$\VSS$ & \texttt{vss} & Value of the Stochastic Solution.\\
\hline
\end{longtable}

Il codice Python deve distinguere chiaramente parametri, dati, variabili
decisionali, output numerici e grafici. Quando possibile, la notazione del
codice deve richiamare direttamente la formulazione matematica presentata nel
manuale.

\section{Convenzioni da evitare}

Per ridurre ambiguit\`{a} e incoerenze, si raccomanda di evitare le seguenti sovrapposizioni:

\begin{enumerate}
\item Non usare $P$ per indicare una probabilit\`{a}. La probabilit\`{a}
\`{e} indicata con $\Prob$, mentre $P$ \`{e} riservato preferibilmente
alle matrici, in particolare alle matrici di transizione.

\item Non usare $S$ per indicare il numero degli scenari quando nel corso
si usa $S_t$ per il prezzo di un'attivit\`{a} finanziaria. L'insieme degli
scenari \`{e} $\mathcal S=\{1,\ldots,N\}$.

\item Non usare $q_s$ per i costi di recourse nella programmazione
stocastica. Il simbolo $q$ \`{e} riservato alla probabilit\`{a}
risk-neutral nel modello binomiale. Per i costi di recourse si usa $d_s$.

\item Non usare $L$ per livelli o limiti quando nella stessa lezione si
lavora con perdite. Per valori soglia, limiti o realizzazioni di perdita
si preferisce $\ell$.

\item Non usare $V_t$ per valori ottimi di problemi di ottimizzazione.
Il simbolo $V_t$ \`{e} riservato al valore temporale di un derivato o di
un processo di valore. Per i valori ottimi si usano $z^*$, $z^{SP}$,
$z^{WS}$ e $z^{EV}$.

\item Non usare $z_s$ per variabili ausiliarie di scenario nei problemi di
CVaR. Per gli eccessi di perdita si usa $u_s$; il simbolo $z$ resta riservato
ai valori ottimi.

\item Non usare $y$ per le variabili duali quando nella stessa parte del corso
si lavora con programmazione stocastica. Le variabili duali sono indicate con
$\lambda$, mentre $y_s$ \`{e} riservato alle decisioni di secondo stadio.

\item Non introdurre simboli locali in conflitto con simboli globali gi\`{a}
definiti in questo documento. Ogni eccezione locale deve essere dichiarata
esplicitamente nella lezione in cui viene usata.
\end{enumerate}

\section{Stato di avanzamento}

\begin{center}
\begin{tabular}{|p{5.0cm}|p{4.0cm}|p{5.0cm}|}
\hline
\textbf{Area} & \textbf{Stato} & \textbf{Note}\\
\hline
Notazione probabilistica & Bozza consolidata & Da validare durante la scrittura delle Lezioni 1--3.\\
\hline
Processi stocastici e filtrazioni & Bozza consolidata & Da coordinare con la prima applicazione Python.\\
\hline
Catene di Markov & Bozza consolidata & Da verificare con gli esempi di rating.\\
\hline
Misure di rischio & Bozza consolidata & Da raccordare con CVaR discreto e formulazione PL.\\
\hline
Pricing binomiale & Bozza consolidata & Da dettagliare nella Lezione 9.\\
\hline
Programmazione lineare e dualit\`{a} & Bozza consolidata & Da adattare agli esempi finanziari scelti.\\
\hline
Programmazione stocastica & Bozza consolidata & Da verificare con la struttura Python finale.\\
\hline
Corrispondenza Python & Bozza iniziale & Da aggiornare quando saranno definiti gli script.\\
\hline
\end{tabular}
\end{center}

\end{document}
